\documentclass[9pt]{extarticle}

\usepackage[utf8]{inputenc}
\usepackage[a4paper, margin=1.2cm]{geometry}
\usepackage[colorlinks=true, urlcolor=blue]{hyperref}
\usepackage{enumitem}
\usepackage{changepage}

\pagestyle{empty}

\setlist[itemize]{label=\textbullet, noitemsep, topsep=0pt}
\setlength{\parindent}{0pt}

% actual document shit

\begin{document}

\begin{center}
	{\huge\bfseries Paul Parisot}
	\\[9pt]
	{\Large Ingénieur Logiciel - DevOps, Cloud \& Backend}
	\\
\end{center}

% section for contact

\par\vspace{-5pt}
\rule{\linewidth}{0.4pt}
\par\vspace{-5pt}

\begin{center}
	Github: \href{https://github.com/paulogarithm}{paulogarithm}
	\\
	\href{mailto:paul.parisot@epitech.eu}{paul.parisot@epitech.eu}
	\\
\end{center}

\par\vspace{-10pt}
\rule{\linewidth}{0.4pt}
\par\vspace{7pt}

% section for the about us

{\large\bfseries À propos}
\par\vspace{-8pt}
\rule{\linewidth}{0.4pt}
\par\vspace{5pt}

Ingénieur logiciel avec de solides bases en Go, Python et C/C++, orienté cloud, conteneurs et automatisation (Docker, Kubernetes, CI/CD, GCP).
\\
Expérience de développement dans des environnements exigeants : contributions open-source sous revue de code et CI/CD stricte (plus de 10 000 lignes de code), démo présentée au CES 2024, enseignement auprès de plus de 400 étudiants.
\\
Formation complémentaire en management à McGill University (comptabilité, économie managériale), un intérêt pour le secteur bancaire et la finance, et une pratique professionnelle de l'anglais (C1). En préparation de l'agrégation d'informatique.
\par\vspace{5pt}

% section for skills

{\large\bfseries Compétences}
\par\vspace{-8pt}
\rule{\linewidth}{0.4pt}
\par\vspace{5pt}

\hspace*{5pt}{\bfseries Maîtrise}: Golang, Python, C/C++, Lua(u), SQL (PostgreSQL), Shell, Haskell, OCaml, Assembleur
\\
\hspace*{5pt}{\bfseries Notions}: Java, COBOL (découverte), AWS (S3)
\\
\hspace*{5pt}{\bfseries Cloud \& DevOps}: GCP, Docker, Kubernetes, Traefik, CI/CD, Git, Environnement Linux
\\
\hspace*{5pt}{\bfseries Langues}: Français, Anglais (C1, professionnel)

\par\vspace{5pt}

% section for the professional experience

{\large\bfseries Expérience Professionnelle}
\par\vspace{-8pt}
\rule{\linewidth}{0.4pt}
\par\vspace{5pt}

{\bfseries Gno.land - Ingénieur Logiciel Junior (Open-Source)}
\begin{itemize}[noitemsep, topsep=0pt]
	\item Développement de 6+ packages core en Gnolang (variante de Go pour smart contracts), incluant printf et une compatibilité avec l'EVM.
	\item Détection et correction de {\bfseries 6 bugs critiques} causant des crashs runtime ; {\bfseries ~10k lignes de code} contribuées.
	\item Rédaction de tests exhaustifs sous des politiques strictes de CI/CD et de revue de code.
\end{itemize}

\par\vspace{5pt}

{\bfseries Faurecia / Forvia - Ingénieur Démo Logicielle}
\begin{itemize}[noitemsep, topsep=0pt]
	\item Conception d'un {\bfseries moteur de calcul d'itinéraire GPS} et d'une interface de recherche de POI, intégrés à une application démo présentée au {\bfseries CES 2024}.
	\item Correction d'un ancien code C++ pour le faire communiquer avec les bons drivers matériels (LED) ; scripts Python de pilotage de hardware.
	\item Livraison de fonctionnalités prêtes pour la production {\bfseries en avance sur le planning}, salué pour la rapidité et la qualité du code.
\end{itemize}

\par\vspace{5pt}

{\bfseries Epitech / IONIS - Assistant Technique \& Formateur}
\begin{itemize}[noitemsep, topsep=0pt]
	\item Enseignement à {\bfseries plus de 400 étudiants} en C, C++, asm et Haskell ; correction et évaluation des rendus.
	\item Création de {\bfseries contenus pédagogiques interactifs} (quiz, cours inversés, tutoriels) et contribution à un {\bfseries taux de réussite historiquement élevé} au sein des promotions coachées par l'ASTEK.
\end{itemize}

\par\vspace{5pt}

% personal projects

{\large\bfseries Projets Personnels}
\par\vspace{-8pt}
\rule{\linewidth}{0.4pt}
\par\vspace{5pt}

\underline{\bfseries Sledge -- Moteur de migration, synchronisation et orchestration multi-cloud}
\par
\begin{adjustwidth}{5pt}{0pt}
    Développement et finalisation d'un outil générique de migration multi-cloud gérant des structures de données complexes et des interactions avec des bases de données. \\
    Architecture et livraison d'une plateforme backend complète en Go, optimisée pour des opérations de flux de données à haut volume, via une couche d'abstraction cloud personnalisée ({\ttfamily provicol}) supportant AWS S3 et GCP. \\
    Utilisation de l'API Docker à l'exécution pour créer et orchestrer dynamiquement des conteneurs d'environnements concurrents.
\end{adjustwidth}
\par
{\bfseries Compétences :} Go, API Docker, Python, SQL, intégrité des données, cloud computing, automatisation système

\par\vspace{8pt}

\underline{\bfseries Area -- Plateforme d'automatisation de services (clone d'IFTTT)}
\par
\begin{adjustwidth}{5pt}{0pt}
    Projet d'équipe : conception et développement seul du backend en Go avec PostgreSQL, en architecture microservices (services en Go et Python), derrière un proxy Traefik. \\
    Intégration d'une quinzaine de services tiers via OAuth, webhooks et requêtes HTTP, y compris des adaptateurs à base de polling planifié (cron) pour les services ne pouvant pas émettre de webhooks.
\end{adjustwidth}
\par
{\bfseries Compétences :} Go, Python, PostgreSQL, microservices, Traefik, OAuth, webhooks, intégration d'API

\par\vspace{8pt}

\underline{\bfseries Zappy -- Serveur TCP multi-clients en C}
\par
\begin{adjustwidth}{5pt}{0pt}
    Serveur TCP mono-thread à boucle événementielle ({\bfseries \texttt{select}}), codé seul : buffer borné par client et mise en quarantaine des clients qui le saturent ; aucune erreur ni fuite sous Valgrind ; tests unitaires et fonctionnels.
\end{adjustwidth}

\par\vspace{8pt}

% education

{\large\bfseries Formation}
\par\vspace{-8pt}
\rule{\linewidth}{0.4pt}
\par\vspace{5pt}

{\bfseries Epitech, Paris} \hfill 2022 -- 2027 \\
Diplôme d'ingénieur en technologies de l'information (diplôme de niveau Master) \\
Certification Niveau 7 - Expert en ingénierie logicielle

\vspace{6pt}

{\bfseries McGill University, Montréal} \hfill 2025 -- 2026 \\
Certificat en Management -- GPA : 3.57/4.0 \\
Cours suivis : Comptabilité, Managerial Economics \& Analysis, Business, Mathématiques, Gestion de projet, Marketing

\end{document}
